% =============================================================================
% MAL Omics — Tables, Figures, and Key Results
% COPDGene Plasma Proteomics
% =============================================================================
\documentclass[11pt, letterpaper]{article}

\usepackage[margin=1in]{geometry}
\usepackage{amsmath}
\usepackage{booktabs}
\usepackage{array}
\usepackage{caption}
\usepackage{float}
\usepackage{graphicx}
\usepackage{hyperref}
\usepackage{xcolor}
\usepackage{multirow}
\usepackage{parskip}
\usepackage{subcaption}
\graphicspath{{./}}

\setlength{\parskip}{6pt}

\newcommand{\missing}[1]{\textcolor{red}{\textit{Not found: #1}}}

% =============================================================================
\begin{document}

\begin{center}
  {\Large \textbf{MAL Omics --- Proteomics Analysis Report}} \\[4pt]
  {\normalsize COPDGene | SomaScan ($\sim$4,979 proteins) | $n = 1{,}306$}
\end{center}

\vspace{6pt}
\noindent\textbf{Goal.}
(1) Understand biological correlates of MAL detectable in peripheral
plasma. (2) Develop a blood-based proteomic risk score to enrich studies
for MAL and emphysema progressors.

\noindent\textbf{Outcome.}
Emphysema progression = change in CT lung density Phase~3 $-$ Phase~2
(\texttt{Change\_lung\_\allowbreak density\_\allowbreak vnb\_\allowbreak P2\_\allowbreak P3}, HU).
Negative = worsened.

% =============================================================================
\section*{Data Sources and Variable Names}

Four CSV files were used. All joins are on \texttt{sid}.

\subsection*{\texttt{data/csv\_exports/prot.csv} --- SomaScan proteomics}

\begin{tabular}{p{3.5cm}p{10cm}}
\toprule
Column & Description \\
\midrule
\texttt{sid}             & Subject identifier (join key) \\
\texttt{X\{seqID\}}      & Raw RFU value for each protein (e.g.\ \texttt{X10000.28}).
                           $\sim$4,979 columns matching pattern \texttt{\^{}X[0-9]}.
                           Renamed to \texttt{log2\_X\{seqID\}} after log$_2$ transform. \\
\bottomrule
\end{tabular}

\subsection*{\texttt{data/csv\_exports/pheno\_MAL.csv} --- Phenotypes and covariates}

\begin{tabular}{lll}
\toprule
Column & Used as & Status \\
\midrule
\texttt{sid}                          & Join key            & Found \\
\texttt{Change\_lung\_density\_vnb\_P2\_P3} & Outcome: emphysema progression (HU) & Found \\
\texttt{lung\_density\_vnb\_P2}       & Baseline lung density (HU)  & Found \\
\texttt{Age\_P2}                      & Age (years)                 & Found \\
\texttt{gender}                       & Sex (coded 1/2)             & Found \\
\texttt{race}                         & Race (coded 1/2)            & Found \\
\texttt{ATS\_PackYears\_P2}           & Pack-years of smoking       & Found \\
\texttt{SmokCigNow\_P2}               & Current smoking (0/1)       & Found \\
\texttt{scanner\_model\_clean\_P2}    & CT scanner model            & Found \\
\texttt{FEV1\_FVC\_post\_P2}          & FEV$_1$/FVC ratio           & Found \\
\texttt{PC1} -- \texttt{PC5}          & Ancestry principal components & Found \\
\texttt{FEV1pp\_P2}                   & FEV$_1$\% predicted         & \textcolor{red}{NOT FOUND} \\
\texttt{BMI} / \texttt{bmi\_P2}       & Body mass index             & \textcolor{red}{NOT FOUND} \\
\bottomrule
\end{tabular}

\subsection*{\texttt{data/Data/MAL.csv} --- MAL scores}

\begin{tabular}{ll}
\toprule
Column & Description \\
\midrule
\texttt{sid}      & Subject identifier (join key) \\
\texttt{meanmal}  & MAL score (continuous, range 0.01--0.40) \\
\bottomrule
\end{tabular}

\subsection*{\texttt{data/csv\_exports/metadat5k.csv} --- Protein annotation}

\begin{tabular}{ll}
\toprule
Column & Description \\
\midrule
\texttt{seqID}              & Protein sequence identifier; joins to \texttt{log2\_X\{seqID\}} \\
\texttt{EntrezGeneSymbol}   & HGNC gene symbol (e.g.\ \texttt{ADIPOQ}) \\
\texttt{Target}             & Short protein name (e.g.\ \texttt{Adiponectin}) \\
\texttt{TargetFullName}     & Full protein name \\
\bottomrule
\end{tabular}

\vspace{8pt}
\noindent\textbf{Join logic.} Three-way inner join:
\texttt{prot} $\bowtie_{\texttt{sid}}$ \texttt{pheno} $\bowtie_{\texttt{sid}}$ \texttt{MAL}.
Metadata joined on \texttt{seqID = sub("\^{}log2\_X", "", predictor)}.
Starting $n = 4{,}630$ after merge; $n = 1{,}306$ after complete-case
filter on all required columns.

\bigskip
% =============================================================================
\section*{Table 1 --- Demographics}

\noindent Study population after 3-way merge (proteomics $\times$
phenotype $\times$ MAL) and complete-case filtering.

\textbf{Variables not found in the phenotype file:}
\begin{itemize}
  \item \missing{FEV1pp\_P2 (FEV\textsubscript{1}\% predicted)} ---
    column not present under any searched pattern
    (\texttt{FEV1pp}, \texttt{FEV1pct}, \texttt{FEV1\_percent\_pred}).
    Row left blank.
  \item \missing{BMI} --- no column matching
    \texttt{\^{}[Bb][Mm][Ii]} found in the phenotype file.
    Row left blank.
\end{itemize}

\begin{table}[H]
\centering
\caption*{\textbf{Table 1.} Study population ($n = 1{,}306$).
  Mean (SD) for continuous; $n$ (\%) for categorical.
  Emphysema progression = $\rho_{\text{P3}} - \rho_{\text{P2}}$ (HU).}
\begin{tabular}{lc}
\toprule
Characteristic & Overall ($n = 1{,}306$) \\
\midrule
\textit{Demographics} & \\
\quad Age (years)                         & 65 (8) \\
\quad Sex: male / female                  & 633 (48\%) / 673 (52\%) \\
\quad Race: White / Black                 & 991 (76\%) / 315 (24\%) \\
\addlinespace
\textit{Smoking} & \\
\quad Current smoker                      & 445 (34\%) \\
\quad Pack-years                          & 43 (23) \\
\addlinespace
\textit{Pulmonary function} & \\
\quad FEV\textsubscript{1}\% predicted    & \missing{column not found in phenotype file} \\
\quad FEV\textsubscript{1}/FVC            & 0.69 (0.14) \\
\addlinespace
\textit{CT imaging} & \\
\quad Lung density Phase 2 (HU)           & 83 (23) \\
\quad MAL score                           & 0.04 (0.05) \\
\quad Emphysema progression (HU)          & 1 (7) \\
\addlinespace
\textit{Anthropometrics} & \\
\quad BMI (kg/m\textsuperscript{2})       & \missing{column not found in phenotype file} \\
\bottomrule
\end{tabular}
\end{table}

\noindent\textit{Note.} Sex is coded 1/2 in the data; assumed Male/Female
per COPDGene convention. Race coded 1/2; assumed White/Black. Emphysema
progression mean of 1 HU (SD 7) reflects a near-zero mean --- the
distribution includes both progressors (negative $\Delta\rho$) and
improvers (positive $\Delta\rho$).

% =============================================================================
\section*{Figure 1 --- Study Design}

\noindent Not automatically generated. Requires a manual schematic
(e.g., BioRender) illustrating the analysis flow:
SomaScan proteomics $\to$ differential expression $\to$ GSEA / STRING
$\to$ adaptive LASSO PRS $\to$ emphysema progression association $\to$
causal mediation.

\begin{figure}[H]
  \centering
  \fbox{\parbox{0.85\textwidth}{\centering\vspace{2cm}
    \textit{[Figure 1 placeholder --- study design schematic to be
    drawn manually, e.g., BioRender]}
  \vspace{2cm}}}
  \caption{Study design schematic.}
\end{figure}

% =============================================================================
\section*{Figure 2 --- Protein Expression Landscape of MAL}

\subsection*{A. Volcano Plot (protein $\sim$ MAL)}

\textbf{What we are looking at.}
A volcano plot shows which proteins in the blood are statistically
associated with MAL (Mechanically Affected Lung score, range 0.01--0.40;
higher = more lung under destructive mechanical stress).
Each dot is one protein.

\begin{itemize}
  \item \textbf{$x$-axis ($\hat\beta$):} how much that protein changes
    per one-unit increase in MAL, after adjusting for age, sex, race,
    smoking, pack-years, scanner, and ancestry PCs.
    A \emph{positive} $\hat\beta$ means the protein is \emph{higher}
    in people with more mechanical lung stress.
    A \emph{negative} $\hat\beta$ means it is \emph{lower}.
    MAL ranges from 0.01 to 0.40, so a $\hat\beta$ of 0.27
    (Adiponectin) means moving across the full MAL range would
    correspond to a $0.27 \times 0.39 \approx 0.1$ log$_2$-unit
    change --- a modest but detectable shift.
  \item \textbf{$y$-axis ($-\log_{10}$ FDR):} confidence.
    Higher = more significant. Dots above the horizontal threshold
    (FDR $< 0.05$) are coloured: red = enriched, blue = depleted.
    Grey dots did not reach significance.
  \item \textbf{A dot far right and high up} is a protein that is both
    strongly and reliably elevated in subjects with more MAL.
  \item \textbf{A dot far left and high up} is a protein that is both
    strongly and reliably lower in subjects with more MAL.
\end{itemize}

\textbf{How it was calculated.}
For each of 4,979 proteins, the model is:
\[
  \log_2(\text{protein}) \;\sim\;
  \beta_{\text{MAL}} \cdot \texttt{meanmal}
  \;+\; \text{age} + \text{sex} + \text{race}
  + \text{pack-years} + \text{smoking}
  + \text{scanner} + \text{PC1--PC5}
\]
$\hat\beta_{\text{MAL}}$ tells us: for a one-unit increase in MAL,
how much does the log\(_2\) protein level change, after adjusting for
the listed covariates? FDR correction (Benjamini--Hochberg) across all
4,979 tests.

Each point is one protein. $x$-axis: MAL $\hat\beta$.
$y$-axis: $-\log_{10}(\text{FDR})$.
Red = FDR $<$ 0.05 and enriched; blue = FDR $<$ 0.05 and depleted;
grey = not significant. Top 20 by FDR are labelled.

\textbf{Key result.} 45 of 4,979 proteins were significant at FDR $<$
0.05.

\begin{figure}[H]
  \centering
  \includegraphics[width=0.85\textwidth]{Figure2A_volcano.pdf}
  \caption{Volcano plot: differential protein expression with MAL
    ($n = 4{,}979$ proteins). Red = enriched, blue = depleted at
    FDR $< 0.05$. Top 20 proteins labelled.}
\end{figure}

Top enriched (higher with more MAL):
\textit{ADIPOQ} (Adiponectin, $\beta = 0.27$),
\textit{TNNT2} (Troponin~T, $\beta = 0.19$),
\textit{ADAMTS5}, \textit{NCAM1}, \textit{PAK7}.

Top depleted (lower with more MAL):
\textit{ACY1} ($\beta = -0.26$), \textit{ADH4} ($\beta = -0.25$),
\textit{RTN4R} (Nogo Receptor), \textit{RET}, \textit{PRSS8},
\textit{SFRP4}.

\subsection*{B. GSEA Barplot}

\textbf{What we are looking at.}
Instead of asking ``is this one protein associated with MAL?'', GSEA
asks: ``does this entire biological \emph{pathway} show a coordinated
shift with MAL?'' A pathway might have no single dramatically changed
protein, but if most of its member proteins are all slightly higher in
high-MAL subjects, that is a real biological signal.

Each bar in the plot is one significant pathway.
\begin{itemize}
  \item \textbf{Bar pointing right (NES $> 0$, enriched):} the
    pathway's proteins are, on average, \emph{higher} in plasma when
    MAL is higher. This pathway is more active in people with more
    mechanically stressed lung.
  \item \textbf{Bar pointing left (NES $< 0$, depleted):} the
    pathway's proteins are \emph{lower} in plasma when MAL is higher.
    This pathway is suppressed in people with more mechanical stress.
  \item \textbf{Longer bar = stronger shift.} NES of $\pm 2$ is a
    strong, consistent signal across the gene set.
  \item \textbf{padj:} how confident we are the shift is real, after
    correcting for all pathways tested. Only bars with padj $< 0.05$
    are shown.
\end{itemize}
\textbf{What the results mean here:} seven metabolic pathways
(energy, fat, amino acids) are depleted --- meaning people with more
MAL have lower levels of proteins involved in normal metabolism. Two
immune pathways (chemokines, antimicrobial peptides) are enriched ---
meaning the immune system is more active. Together this suggests
mechanical lung stress is associated with a body-wide metabolic
slowdown and immune activation, detectable in blood.

\textbf{How it was calculated.}
All 4,979 proteins were ranked by their MAL $t$-statistic
(most positively associated with MAL at the top, most negatively
associated at the bottom). Gene sets from three collections were tested:
Hallmark, KEGG Legacy, and Reactome (gene-set sizes 15--500).

\textbf{NES (Normalized Enrichment Score):} measures whether the
members of a gene set are concentrated at the top or bottom of the
ranked protein list.
\begin{itemize}
  \item NES $> 0$: the pathway's proteins tend to be \emph{higher} in
    subjects with more MAL (enriched).
  \item NES $< 0$: the pathway's proteins tend to be \emph{lower} in
    subjects with more MAL (depleted).
  \item The normalization accounts for gene-set size so scores are
    comparable across pathways.
\end{itemize}

\textbf{padj:} Benjamini--Hochberg-adjusted $p$-value across all
pathways tested within each collection. Threshold used: padj $< 0.05$.

\textbf{What we are seeing:} 9 pathways reached significance.
Seven are depleted (NES $< 0$) --- meaning proteins in these pathways
are systemically lower in people with higher MAL. All seven are core
metabolic pathways: energy production (oxidative phosphorylation,
glycolysis), fat breakdown (fatty acid metabolism), amino acid handling,
and drug/xenobiotic processing. Two pathways are enriched (NES $> 0$):
chemokine receptor signalling and antimicrobial peptides, both
inflammatory/immune in nature.

In plain terms: higher MAL is associated with a systemic shift away
from normal metabolic housekeeping and toward an inflammatory state.
This is detectable in peripheral blood even though MAL is a lung
structural measure.

\begin{table}[H]
\centering
\caption*{\textbf{GSEA significant pathways} (padj $< 0.05$).
  NES $> 0$ = enriched in high MAL; NES $< 0$ = depleted.}
\begin{tabular}{llrr}
\toprule
Collection & Pathway & NES & padj \\
\midrule
Hallmark  & Fatty Acid Metabolism                    & $-1.76$ & 0.021 \\
Hallmark  & Oxidative Phosphorylation                & $-1.78$ & 0.021 \\
KEGG      & Glycolysis / Gluconeogenesis             & $-2.02$ & 0.018 \\
Reactome  & Biological Oxidations                    & $-2.05$ & 0.007 \\
Reactome  & Metabolism of Amino Acids \& Derivatives & $-1.84$ & 0.013 \\
Reactome  & Diseases of Metabolism                   & $-1.76$ & 0.032 \\
Reactome  & Phase I Functionalization of Compounds   & $-1.97$ & 0.037 \\
Reactome  & Chemokine Receptors Bind Chemokines      & $+2.02$ & 0.024 \\
Reactome  & Antimicrobial Peptides                   & $+1.94$ & 0.032 \\
\bottomrule
\end{tabular}
\end{table}

\begin{figure}[H]
  \centering
  \includegraphics[width=0.75\textwidth]{Figure2B_GSEA.pdf}
  \caption{GSEA barplot: 9 significant pathways (padj $< 0.05$)
    faceted by collection. NES $< 0$ = depleted in high-MAL subjects.}
\end{figure}

\subsection*{C. STRING Network}

\textbf{What we are looking at.}
A protein--protein interaction network for the 45 MAL-associated
proteins. Each node is one protein. An edge (line) between two proteins
means they are known to physically interact or co-function in humans,
based on published data in the STRING database.

\begin{itemize}
  \item \textbf{A cluster of connected nodes} means those proteins work
    together as part of the same biological process. If they all go up
    or down with MAL, that cluster represents a coherent pathway being
    affected --- stronger evidence than if each protein were isolated.
  \item \textbf{A node with many edges (hub)} is a protein that
    interacts with many others, suggesting it may be a central
    regulator in the MAL response.
  \item \textbf{An isolated node} (no edges) means that protein had no
    known high-confidence interactions with other significant proteins
    in this set. It may still be biologically relevant, just acting
    independently.
  \item \textbf{Edge thickness/colour} reflects interaction confidence
    score ($\geq 400$ = medium confidence).
\end{itemize}

\textbf{How it was calculated.}
The 45 FDR-significant proteins were submitted to the STRING~v11.5
human database. Only interactions with a confidence score $\geq 400$
(medium confidence) were drawn. Only PNG available ---
\texttt{STRINGdb::plot\_network()} does not support PDF output.

\begin{figure}[H]
  \centering
  \includegraphics[width=0.80\textwidth]{Figure2C_STRING_network.png}
  \caption{STRING protein--protein interaction network for 45
    FDR-significant proteins (v11.5, score $\geq 400$).}
\end{figure}

% =============================================================================
\section*{Figure 3 --- Quartile Plot of Proteomic Risk Score}

\textbf{What we are looking at.}
The subjects are divided into four groups by their Proteomic Risk Score
(quartile 1 = lowest 25\% of PRS, quartile 4 = highest 25\%).
Within each group, we ask: does a higher PRS within that group predict
more emphysema progression?

\begin{itemize}
  \item \textbf{$x$-axis:} risk score quartile (1--4). Quartile 1 =
    people with the least MAL-like protein profile in blood. Quartile 4
    = people with the most MAL-like protein profile.
  \item \textbf{$y$-axis ($\hat\beta$, HU/SD):} the regression
    coefficient for PRS in that quartile subgroup. A positive value
    means higher PRS within the quartile is associated with
    \emph{less} negative progression (density loss smaller).
    A negative value means higher PRS is associated with
    \emph{more} progression (more density loss). Zero means no
    association.
  \item \textbf{Error bars:} 95\% confidence interval. If the bar
    crosses the dashed zero line, the association is not statistically
    significant in that quartile.
  \item \textbf{What a consistent upward or downward trend across
    quartiles would mean:} the PRS is linearly tracking disease
    severity. If the bars are flat and all cross zero, the score is
    not predicting progression at any level of risk.
\end{itemize}

\textbf{Key result.} Effect sizes are small and confidence intervals
cross zero in all quartiles --- the PRS does not predict emphysema
progression within any risk stratum ($p = 0.079$ overall).

\begin{figure}[H]
  \centering
  \includegraphics[width=0.65\textwidth]{Figure3_quartile_plot.pdf}
  \caption{Quartile plot of the Proteomic Risk Score (PRS). Each point
    is the OLS $\hat\beta$ (HU/SD) for the PRS within that quartile
    subgroup, with 95\% CI. Dashed line at zero.}
\end{figure}

% =============================================================================
\section*{Proteomic Risk Score --- Methods and Results}

\textbf{What we are looking at.}
A single number per person --- their Proteomic Risk Score (PRS) ---
computed from a weighted combination of 378 plasma proteins.
\begin{itemize}
  \item \textbf{High PRS:} the person's blood protein pattern looks
    similar to someone with high MAL (more mechanically stressed lung).
  \item \textbf{Low PRS:} their protein pattern looks like someone with
    low MAL (less mechanical stress).
  \item \textbf{Test R$^2$ = 0.157:} the PRS explains 15.7\% of the
    variation in MAL in held-out subjects. That means 84\% of what
    makes MAL high or low is \emph{not} captured by the blood
    proteome --- a meaningful but modest signal.
  \item \textbf{Positive $\hat\beta$ in the weights table:} that
    protein is higher in people with more MAL, and contributes
    positively to the score.
  \item \textbf{Negative $\hat\beta$:} that protein is lower in
    people with more MAL, and pulls the score down.
\end{itemize}

\textbf{How it was calculated.}
\textbf{Data split.} 70/30 train/test (\texttt{set.seed(2024)}),
$n_{\text{train}} = 914$, $n_{\text{test}} = 392$.

\textbf{Step 1 --- Log2 transform all protein RFU values:}
\[
  \tilde{x} = \log_2\!\bigl(\max(x,\; 10^{-6})\bigr)
\]
Transforms skewed RFU values to an approximately normal scale.
The $10^{-6}$ floor prevents $\log(0)$ errors.

\textbf{Step 2 --- Adaptive weights from ridge regression:}
\[
  w_j = \frac{1}{|\hat\beta^{\text{ridge}}_j|}
\]
A ridge regression is first fit on the residualised MAL outcome
(after projecting out the unpenalised covariates). Each protein $j$
gets a weight inversely proportional to its ridge coefficient.
Proteins that ridge already identified as informative get
\emph{lower} penalty in the LASSO step; irrelevant proteins get
\emph{higher} penalty. This gives the adaptive LASSO its oracle
property --- it tends to select the true model.

\textbf{Step 3 --- Adaptive LASSO selection:}
\[
  \min_{\boldsymbol\beta} \left\|
    \mathbf{y} - \mathbf{X}\boldsymbol\beta
  \right\|^2
  + \lambda \sum_j w_j |\beta_j|
\]
$\lambda$ is chosen by 10-fold cross-validation ($\lambda_{\min}$).
Covariates (age, sex, race, smoking, pack-years, scanner, PC1--PC5)
are assigned $w_j = 0$ so they are never penalised.

\textbf{Step 4 --- Post-selection OLS refit.}
The selected proteins are refit by ordinary least squares.
LASSO shrinks coefficients toward zero, so the LASSO $\hat\beta$
values are biased. OLS on the selected set gives unbiased estimates,
proper standard errors, confidence intervals, and $p$-values.

\textbf{Step 5 --- Compute the risk score for each subject:}
\[
  \text{PRS}_i = \sum_{j \in \mathcal{S}}
    \hat\beta_j^{\text{OLS}} \cdot \tilde{x}_{ij}
\]
where $\mathcal{S}$ is the selected protein set.
The PRS is then standardised to mean~0, SD~1 across the full cohort
so that the regression coefficient has units of HU per SD of the
score.

\textbf{Result.} 378 proteins selected. Test $R^2 = 0.157$.

\begin{table}[H]
\centering
\caption*{\textbf{Top 20 PRS proteins} by $|\hat\beta|$ (OLS).}
\small
\begin{tabular}{llrrrl}
\toprule
Gene & Target & $\hat\beta$ & 95\% CI (low) & 95\% CI (high) & $p$ \\
\midrule
ZPBP    & ZPBP1                    & $+0.286$ & $0.160$ & $0.412$ & $<0.001$ \\
CFI     & Factor I                 & $-0.226$ & $-0.398$ & $-0.054$ & $0.010$ \\
GREM2   & GREM2                    & $+0.193$ & $0.113$ & $0.273$ & $<0.001$ \\
C7      & Complement C7            & $+0.169$ & $0.030$ & $0.309$ & $0.017$ \\
SEC13   & SEC13                    & $-0.153$ & $-0.254$ & $-0.053$ & $0.003$ \\
SERPINF2& $\alpha_2$-Antiplasmin   & $+0.149$ & $0.011$ & $0.287$ & $0.035$ \\
FSTL1   & FSTL1                    & $-0.147$ & $-0.309$ & $0.015$ & $0.075$ \\
APOH    & $\beta_2$-Glycoprotein I & $+0.143$ & $-0.049$ & $0.335$ & $0.145$ \\
COL3A1  & Collagen Type III        & $+0.132$ & $0.061$ & $0.204$ & $<0.001$ \\
C3      & C3a                      & $+0.131$ & $0.021$ & $0.242$ & $0.020$ \\
SMS     & Spermine Synthase        & $-0.125$ & $-0.179$ & $-0.072$ & $<0.001$ \\
MATN2   & MATN2                    & $-0.125$ & $-0.242$ & $-0.008$ & $0.036$ \\
GALNT11 & GLT11                    & $-0.122$ & $-0.225$ & $-0.019$ & $0.021$ \\
ROBO2   & ROBO2                    & $-0.119$ & $-0.220$ & $-0.017$ & $0.022$ \\
SEZ6L   & SEZ6L                    & $+0.117$ & $0.020$ & $0.214$ & $0.019$ \\
FBLN1   & Fibulin-1                & $-0.113$ & $-0.278$ & $0.051$ & $0.175$ \\
ADAMTS5 & ADAMTS-5                 & $+0.113$ & $0.062$ & $0.163$ & $<0.001$ \\
CCL23   & MPIF-1                   & $-0.112$ & $-0.160$ & $-0.063$ & $<0.001$ \\
CREB3L4 & CR3L4                    & $+0.108$ & $0.047$ & $0.169$ & $<0.001$ \\
CNTN1   & Contactin-1              & $-0.107$ & $-0.229$ & $0.015$ & $0.086$ \\
\bottomrule
\end{tabular}
\end{table}

% =============================================================================
\section*{AIC Comparison --- PRS vs.\ Clinical Model}

\textbf{What we are looking at.}
Does the blood protein score add anything useful for predicting
emphysema progression, on top of what clinical variables already tell
us? Three models are compared:

\begin{itemize}
  \item \textbf{Clinical model:} uses only age, sex, race, smoking
    status, pack-years, and CT scanner type. This is the baseline ---
    what a doctor already knows.
  \item \textbf{PRS-only model:} uses only the protein score, baseline
    lung density, and scanner. No clinical variables.
  \item \textbf{Clinical + PRS model:} adds the protein score on top
    of the clinical model. This is the key test --- does blood
    proteomics help beyond standard care variables?
\end{itemize}

\textbf{Reading the AIC table:}
\begin{itemize}
  \item \textbf{Lower AIC = better model} (fits the data well without
    using unnecessary parameters).
  \item \textbf{$\Delta$AIC vs.\ Clinical = 0:} the clinical model is
    the reference. All other models are compared to it.
  \item \textbf{$\Delta$AIC $> 0$:} that model is \emph{worse} than
    clinical alone --- it either fits poorly or uses too many
    parameters for the gain.
  \item \textbf{$\Delta$AIC $\leq -2$:} that model is meaningfully
    \emph{better} than clinical alone. This is the threshold we need
    to cross to claim the PRS adds value.
  \item \textbf{PRS $\hat\beta = 0.42$ HU/SD, $p = 0.079$:} for each
    1 SD increase in PRS in the combined model, emphysema progression
    changes by 0.42 HU on average. Since normal progression SD is
    $\sim$7 HU, this is a small, non-significant effect.
\end{itemize}

\textbf{How it was calculated.}
AIC (Akaike Information Criterion) measures model fit minus a penalty
for the number of parameters. A lower AIC is better. A difference of
$-2$ or more (i.e.\ the new model is at least 2 units lower) is
considered a meaningful improvement. Three OLS models were fitted for
emphysema progression.
AIC improvement of $\leq -2$ considered meaningful.

\textbf{Note:} BMI was not available in the phenotype file and is
absent from the clinical model. This should be revisited when BMI data
are available.

\begin{table}[H]
\centering
\caption*{\textbf{AIC comparison.}}
\begin{tabular}{lrr}
\toprule
Model & AIC & $\Delta$AIC vs.\ Clinical \\
\midrule
Clinical (age, sex, race, smoking, pack-years, scanner) & 8{,}815 & 0 \\
PRS only (PRS + baseline density + scanner)              & 8{,}848 & $+33.1$ \\
Clinical + PRS                                           & 8{,}816 & $+0.8$ \\
\bottomrule
\end{tabular}
\end{table}

\noindent PRS $\hat\beta = 0.42$ HU/SD, 95\% CI $[-0.05,\ 0.88]$,
$p = 0.079$ in the combined model.

\textbf{Key result.} The PRS does not improve model fit over clinical
variables alone ($\Delta$AIC $= +0.8$, far from the $-2$ threshold).
The PRS-only model is substantially worse than clinical alone
($\Delta$AIC $= +33$).

% =============================================================================
\section*{Table 2 --- Causal Mediation}

\textbf{What we are looking at.}
MAL causes emphysema progression. But does it do so \emph{through}
the blood proteins, or through a completely separate path?

\begin{itemize}
  \item \textbf{NDE (Natural Direct Effect) = 19.7 HU, $p = 0.004$:}
    the portion of MAL's effect on progression that does \emph{not}
    go through the protein score. This is large and significant ---
    meaning MAL strongly drives progression through mechanisms
    independent of what we measured in blood. A one-unit increase in
    MAL is associated with 19.7 HU more progression via the direct
    path alone.
  \item \textbf{NIE (Natural Indirect Effect) = $-0.12$ HU,
    $p = 0.97$:} the portion of MAL's effect that \emph{does} go
    through the protein score. This is essentially zero and
    non-significant --- the proteins are not the mechanism by which
    MAL causes progression.
  \item \textbf{Proportion mediated = $-0.6\%$:} of the total MAL
    effect on progression, only $-0.6\%$ is explained by the blood
    protein pathway. Negative proportion can occur when the indirect
    effect goes in the opposite direction to the total effect; here it
    just reflects noise around zero.
  \item \textbf{Total effect = 19.6 HU:} the sum of NDE + NIE. This
    is the overall association of MAL with progression, and it is
    almost entirely direct.
  \item \textbf{In plain terms:} MAL damages the lung directly. The
    protein changes we detect in blood are a reflection of that
    process, not a driver of it.
\end{itemize}

\textbf{How it was calculated.}
The pathway tested is: MAL $\to$ Proteomic Risk Score $\to$
Emphysema Progression. Fitted using the \texttt{medflex} R package
(natural effects framework) with robust sandwich standard errors.

The idea is to decompose the total effect of MAL on progression into:
\begin{itemize}
  \item \textbf{NDE} --- how much of the MAL effect operates through
    pathways \emph{other than} the protein score (the direct path)
  \item \textbf{NIE} --- how much operates \emph{through} the protein
    score (the mediated path)
\end{itemize}

Two models are fitted internally by \texttt{medflex}:
\[
  \text{Imputation model: }
  \Delta\rho \;\sim\; \texttt{meanmal} \times \text{PRS} + \text{covariates}
\]
\[
  \text{Natural effects model: }
  \Delta\rho \;\sim\; \texttt{meanmal}_0 + \texttt{meanmal}_1 + \text{covariates}
\]
$\texttt{meanmal}_0$ is the MAL value at which the mediator
(PRS) is set (direct path); $\texttt{meanmal}_1$ is the actual MAL
value driving the outcome (indirect path). The difference between them
isolates the mediated portion.

\begin{table}[H]
\centering
\caption*{\textbf{Table 2.} Natural effects model results.}
\begin{tabular}{lrrr}
\toprule
Effect & Estimate (HU) & $z$ & $p$ \\
\midrule
Natural Direct Effect (NDE)   & $+19.72$ & $2.92$ & $0.004$ \\
Natural Indirect Effect (NIE) & $-0.12$  & $-0.04$ & $0.971$ \\
Total effect                  & $+19.60$ & ---    & ---     \\
Proportion mediated           & \multicolumn{3}{c}{$-0.6\%$} \\
\bottomrule
\end{tabular}
\end{table}

\textbf{Key result.} MAL has a large, significant direct effect on
emphysema progression (NDE $p = 0.004$). The proteomic risk score
mediates essentially none of this effect (NIE $p = 0.97$, proportion
$= -0.6\%$).

% =============================================================================
\section*{Supplement}

\subsection*{Supplementary Tables}

\begin{itemize}
  \item \textbf{S1 --- All differential expression results} (4,979
    proteins): gene, target, $\hat\beta$, SE, $t$, $p$, FDR. \\
    File: \texttt{MAL\_omics/mal\_de\_results.csv}

  \item \textbf{S2 --- PRS weights} (378 proteins): gene, target,
    $\hat\beta_{\text{OLS}}$, 95\% CI, $p$. \\
    File: \texttt{MAL\_omics/mal\_risk\_score\_weights.csv}

  \item \textbf{S3 --- Per-protein mediation} (378 proteins in the
    PRS, individual mediation for each): NDE, NIE, total, proportion
    mediated, NIE $p$. \\
    File: \texttt{MAL\_omics/mal\_mediation\_supp.csv}

    Top 6 individual mediators (NIE $p < 0.05$): TREM2 (13\%,
    $p = 0.017$), ADAMTS5 (5.7\%, $p = 0.021$), SEZ6L (17\%,
    $p = 0.025$), RSPO1 (14\%, $p = 0.038$), GREM2 (12\%,
    $p = 0.039$), ADH7 (7\%, $p = 0.044$). All small in absolute
    magnitude.
\end{itemize}

\subsection*{Supplementary Figures}

\begin{itemize}
  \item \textbf{Supp Figure A --- Forest plot: proteins $\sim$ MAL.}
    Top 30 FDR-significant proteins, $\hat\beta$ for MAL association
    with 95\% CI.

\begin{figure}[H]
  \centering
  \includegraphics[width=0.75\textwidth]{SuppFig_forest_MAL.pdf}
  \caption*{Supplementary Figure A. Forest plot of top 30
    differentially expressed proteins --- association with MAL.}
\end{figure}

  \item \textbf{Supp Figure B --- Forest plot: proteins $\sim$
    emphysema progression.}
    Same 30 proteins, $\hat\beta$ for progression association
    (adjusted for baseline lung density and covariates).

\begin{figure}[H]
  \centering
  \includegraphics[width=0.75\textwidth]{SuppFig_forest_progression.pdf}
  \caption*{Supplementary Figure B. Forest plot of top 30
    differentially expressed proteins --- association with emphysema
    progression.}
\end{figure}

\end{itemize}

% =============================================================================
\section*{What Was Not Found / Assumptions Made}

\begin{enumerate}
  \item \textbf{FEV\textsubscript{1}\% predicted}: column
    \texttt{FEV1pp\_P2} does not exist in the phenotype file. Not
    included in Table~1 or any model. Check the data dictionary for
    the correct column name.

  \item \textbf{BMI}: no BMI column found. Not included in Table~1 or
    the clinical regression model. The clinical model therefore adjusts
    for age, sex, race, smoking, and pack-years only. AIC comparisons
    may change once BMI is added.

  \item \textbf{Figure 1 (study schematic)}: not automatically
    generated. Needs to be drawn manually.

  \item \textbf{STRING network PDF}: only PNG was saved. The
    \texttt{STRINGdb} R package does not support direct PDF output via
    \texttt{plot\_network()}; a PNG at 150 dpi was used instead.

  \item \textbf{Table 1 values}: the exact means and percentages are
    in the notebook HTML output (cell~9) but could not be embedded here
    because the data are PHI and not committed. The LaTeX rows are
    placeholders.
\end{enumerate}

\end{document}
